\scanpage{23}\mspage{22}
It follows that $\overline{\{\deleted{\unsure{b}}a\}}=\overline{\{b\}}\Longleftrightarrow a=b$ (thus the topological space $\mathcal L$ satisfies the axiom $T_0$, it is ``presober'' i.e. injects into the presobrification) -- I do not know whether this is
the case for every segment piece, nor whether a segment piece, with its canonical topology, is sober (thus, moreover, every closed \uline{irreducible} subset would have a generic point \dots)
\figure{\begin{tikzpicture}[x=1cm,y=1cm]\draw (0,0)--(1.25,-.03)--(1.56,1)--(3.4,1);\draw (1.18,-.01)--(1.49,1);\end{tikzpicture}}
\margin{\unsure{To|To} be checked,}
the ``segment'' $[a,b]$ are not necessarily closed for the canonical topology$^{\dagger}$\unsure{,|.}
It is only when $\mathcal L$ is \uline{totally} ordered that one can conclude that the points are closed (converse ?) and in fact, $\mathcal L$ is $\unsure{=|\simeq}$ \uline{Hausdorff}, and moreover the $[a,b]$, $\mathcal L_{\geq a}$, $\mathcal L_{\leq\deleted{\unsure{b}}a}$ are closed.

One can define another topology, having as a fundamental basis of closed sets the $[a,b]$, $\mathcal L_{\geq a}$, $\mathcal L_{\leq a}$ (of which one takes arbitrary intersections of finite unions, to obtain all the closed sets)

% pale vertical watermark
% the paper's watermark (EXTRA STRONG), not his text
